\documentclass[12pt,a4paper]{article}
\usepackage[utf8]{inputenc}
\usepackage[T1]{fontenc}
\usepackage{amsmath}
\usepackage{amssymb}
\usepackage{graphicx}
\usepackage[spanish]{babel}
\usepackage[a4paper, margin=2.5cm]{geometry}
\usepackage[european]{circuitikz}
\usetikzlibrary{babel}

\title{Análisis de Circuito: Equivalente de Thévenin entre A y B}
\author{Gemini Code Assist}
\date{\today}

\begin{document}

\maketitle

\section{Descripción del Circuito}

El circuito está compuesto por un nodo de entrada A y un nodo de salida B. Entre ellos se encuentra un nodo intermedio C. El objetivo es determinar la tensión de Thévenin ($E_{th2}$) y la resistencia equivalente ($R_{eq2}$) vistas desde los terminales A y B.

\begin{center}
    \begin{circuitikz}[scale=1.2, transform shape]
        % Nodos A y B a la izquierda
        \node[circ, label=left:A] (A) at (0, 2) {};
        \node[circ, label=left:B] (B) at (0, -2) {};
        
        % Resto del circuito a la derecha
        \draw (A) to[R, l=$R_3$] (3,2) node[circ, label=above:C] (C) {};
        \draw (C) to[R, l=$R_5$] (6,2) to (6,-2) -- (B);
        \draw (C) to[R, l=$R_2$] (3,0) to[V, v<=$E_2$] (6,0) to (6,-2);
    \end{circuitikz}
\end{center}

\section{Cálculo de la Tensión de Thévenin ($E_{th2}$)}

La tensión de Thévenin, $E_{th2}$, es la tensión en circuito abierto entre los nodos A y B, es decir, $V_{AB} = V_A - V_B$.

Como no hay ninguna fuente de corriente o tensión en la rama que contiene a $R_3$, no circula corriente por ella en condiciones de circuito abierto. Por lo tanto, no hay caída de tensión en $R_3$, y el potencial en el nodo A es el mismo que en el nodo C ($V_A = V_C$).

Entonces, la tensión de Thévenin es:
\begin{equation}
    E_{th2} = V_{AB} = V_C - V_B = V_{CB}
\end{equation}

La tensión $V_{CB}$ es la tensión a través de las dos ramas en paralelo que conectan C y B. Ambas ramas tienen la misma tensión. Podemos calcularla usando la fórmula de Millman entre los nodos C y B, o simplemente analizando el divisor de tensión en la rama que contiene la fuente.

Considerando la rama con $E_2$, $R_2$ y $R_5$ como un circuito cerrado (ya que $R_5$ está en paralelo), la tensión $V_{CB}$ es la tensión que cae en la resistencia $R_5$. Las resistencias $R_2$ y $R_5$ forman un divisor de tensión con la fuente $E_2$.

\begin{equation}
    E_{th2} = V_{CB} = E_2 \cdot \frac{R_5}{R_2 + R_5}
\end{equation}

\section{Cálculo de la Resistencia Equivalente ($R_{eq2}$)}

Para calcular la resistencia equivalente ($R_{eq2}$), también llamada resistencia de Thévenin ($R_{Th2}$), apagamos todas las fuentes independientes. La fuente de tensión $E_2$ se reemplaza por un cortocircuito. Luego, calculamos la resistencia total vista desde los terminales A y B.

\begin{center}
    \begin{circuitikz}[scale=1.2, transform shape]
        % Nodos A y B a la izquierda
        \node[circ, label=left:A] (A) at (0, 2) {};
        \node[circ, label=left:B] (B) at (0, -2) {};
        
        % Resto del circuito a la derecha (con fuente apagada)
        \draw (A) to[R, l=$R_3$] (3,2) node[circ, label=above:C] (C) {};
        \draw (C) to[R, l=$R_5$] (6,2) to (6,-2) -- (B);
        \draw (C) to[R, l=$R_2$] (3,0) to[short] (6,0) to (6,-2);
    \end{circuitikz}
\end{center}

Desde el punto de vista de los terminales A y B, la resistencia $R_3$ está en serie con la combinación de resistencias que se ve desde el nodo C hacia el nodo B.

Desde el nodo C, las resistencias $R_2$ y $R_5$ están conectadas en paralelo entre C y B. Su resistencia equivalente es:
\begin{equation}
    R_{2 \parallel 5} = \frac{R_2 \cdot R_5}{R_2 + R_5}
\end{equation}

Finalmente, la resistencia total equivalente $R_{eq2}$ es la suma de $R_3$ y esta combinación en paralelo:
\begin{equation}
    R_{eq2} = R_3 + (R_2 \parallel R_5) = R_3 + \frac{R_2 \cdot R_5}{R_2 + R_5}
\end{equation}

\end{document}